\section{Introduction}
\label{sec:intro}

An assessment item carries a tag that names the competency it is said to measure: a standard code, a cognitive domain, or a rubric criterion. A correct answer is then read as evidence of that competency. Accountability systems report proficiency by standard, districts buy materials described as standards-aligned and take the aligned assessment as evidence that the materials teach the standard, and research uses a tagged assessment as the outcome measure of an intervention; each use assumes that passing requires the competency. The testing standards treat that reading as a claim to be validated for the intended use~\cite{aera2014standards}, and validity theory treats it as an argument that can fail even when the tag is sincere~\cite{Kane_1992,Kane_2013,Messick_1995}: the item may be answerable from a surface cue, from the options without the stem, or from a routine that is never bound to the situation the tag names.

We test the assumption directly and without student data by searching for \emph{assessment bypasses}: solvers that lack the tagged competency by an explicit restriction on the operations they may perform or the information they receive, never see the published key, and still pass under the published scoring rule, so that a pass on those items does not require the part of the competency the solver was denied. Neither restriction kind belongs to a subject: a lexical policy that answers a causal-inference item by locating the phrase random assignment lacks the operation of causal identification, and a reader given a reading-comprehension question and its options with the passage withheld lacks the information the item is about; partial-input baselines~\cite{kaushik-lipton-2018-much} already show that such readers beat chance. The claim is one-sided: a bypass says nothing about what students learned, and a search that finds none does not certify the tag beyond the solvers that were tried. Alignment tools, which match item text to a standard, ask whether the item mentions the topic or offers an opportunity to use the competency; psychometric methods ask whether performance discriminates students who have the competency from those who do not, and need student responses. Whether passing requires the competency sits between those questions, needs no students, and is what a bypass refutes.

We run the search first on public-release mathematics items, because several authorities publish standards, keys, scoring rules, and item text in machine-checkable form, and the catalog of shortcut strategies for those items is developed enough to write capability restrictions as small programs. We find bypasses. On New York Regents Algebra~I and Algebra~II, a language model shown the question with every given quantity replaced by a placeholder recovers the published key well above chance and above the most common key letter, in two model families, and the recovery survives replacing every option value with a different value of the same form. On those items a pass does not require computing on the given quantities; whether it requires the tagged operation is a question the item codings leave open. Small programs that lack the tagged operation match keys on a few cells: a rule that picks the second-smallest numeric option on two cells, and substitution of the options into the equation read from the page image on a small subset of Algebra~I. The cognitive-domain cells of the Trends in International Mathematics and Science Study (TIMSS)~\cite{iea2009timss2011frameworks}, at the 50 to 80 items per domain that are public, yield no bypass. Every claimed rate is measured on items re-read from the printed page, because the text layer of a released PDF often differs from the print, and the language-model rates only where a second transcription of the page agrees with the first.

Our contributions are the object and the bar, which are general, and the census and its findings, which are the mathematics instantiation.

\begin{itemize}
\item \textbf{The assessment bypass.} A one-sided, student-free test of any tag: a solver that lacks the tagged competency by an explicit capability or information restriction and passes under the published scoring rule refutes the claim that passing requires what it was denied. Validity theory states the argument, and cognitive diagnosis and test-wiseness research estimate students or list cues; none scores this object (Section~\ref{sec:related}).
\item \textbf{A pre-registered bar.} A bypass must fire on at least 10 items with a lower 95\% confidence bound above chance, programs are hashed before scoring, the three solver families are Holm-corrected separately, and every claimed row is measured on the item as printed (Section~\ref{sec:method}).
\item \textbf{A census.} The first instantiation of the method: 2{,}405 keyed public-release mathematics items from nine authorities, each tag quoted from its authority, scored by item-writing cues, test-taking strategies, option substitution into the printed equation, and language models denied the given quantities.
\item \textbf{Findings.} Language-model bypasses on Algebra~I and Algebra~II that survive a change of every option value, two catalog rows awaiting replication, one strategy row that failed a powered replication, and a TIMSS null at public $n$ (Section~\ref{sec:results}).
\end{itemize}
